\documentclass[11pt]{article}

\usepackage[margin=1in]{geometry}
\usepackage[T1]{fontenc}
\usepackage[utf8]{inputenc}
\usepackage{lmodern}
\usepackage{amsmath, amssymb}
\usepackage{booktabs}
\usepackage{threeparttable}
\usepackage{caption}
\usepackage{array}
\usepackage{graphicx}
\usepackage{natbib}
\usepackage{hyperref}

\newcommand{\TableNote}[1]{%
    \par\vspace{0.35em}
    \begin{minipage}{0.96\textwidth}
    \footnotesize\emph{Notes.} #1
    \end{minipage}
}

\hypersetup{
    colorlinks=true,
    linkcolor=blue,
    citecolor=blue,
    urlcolor=blue
}

\title{Robots.txt Restrictions, Filing-Window Illiquidity, and Drift}
\author{}
\date{}

\begin{document}
\maketitle

\section{Robots.txt Restrictions, Filing-Window Illiquidity, and Drift}
\label{sec:robots_txt_illiquidity_drift}

\subsection{Treatment Definitions}
\label{subsec:robots_treatment}

I construct treatment directly from archived \texttt{robots.txt} rules.
For each firm-month, I identify whether at least one observed corporate
host restricts crawler access through a \texttt{Disallow} rule.  The
baseline treatment family is intentionally about search-bot access in
general, rather than a single AI product.  The main variable is
\textit{SearchBotSpec}: it turns on when a firm adopts a bot-specific
\texttt{User-agent} group that blocks at least one search-oriented crawler.
I require an observed clean-to-block transition: treated firms must have
at least one valid pre-treatment month in which the relevant crawler group
is not blocked.  Firms already blocked at the beginning of the panel are
excluded as left-censored.  Firms with valid robots.txt observations that
never block the relevant crawler group form the clean-never control group.

For firm \(i\), let \(G_i^{SB}\) denote the first month in which a
search-bot-specific block is observed.  For filing event \(e\) in month
\(m(e)\), the treatment indicator is
\[
    \text{SearchBotSpec}_{i,e}
    =
    \mathbf{1}\{m(e)>G_i^{SB}\}.
\]
The strict inequality is deliberate: the first blocking month itself is
not treated as post.  Filing-window regressions also exclude filings in
the first blocking month.

Table~\ref{tab:robots_treatment_definitions} defines the treatment
variables used below.  The key distinction is between \textit{SearchBotSpec}
and \textit{BotSpec}.  \textit{SearchBotSpec} is limited to the search-bot
side of the crawler taxonomy, while \textit{BotSpec} uses a broader
tracked crawler set.  Both require bot-specific \texttt{User-agent} rules.
By contrast, the wildcard variables isolate restrictions inherited from
\texttt{User-agent: *}; these are broader site policies rather than
bot-specific targeting.

\begin{table}[!htbp]
\centering
\caption{Robots.txt Treatment Definitions}
\label{tab:robots_treatment_definitions}
\begin{threeparttable}
\small
\begin{tabular}{llp{8.1cm}}
\toprule
Variable & Label & Definition \\
\midrule
\textit{SearchBotSpec} &
Search-bot-specific block &
Firm adopts a bot-specific \texttt{User-agent} rule that blocks at least
one search-oriented crawler.  This is the main general search-bot
treatment. \\
\textit{BotSpec} &
Broader bot-specific block &
Firm adopts a bot-specific \texttt{User-agent} rule that blocks at least
one crawler in the broader tracked crawler set.  This is a broader
bot-specific robustness variable and is not limited to the search-bot side. \\
\textit{RobotsAny} &
Any robots.txt restriction &
Firm blocks at least one tracked crawler through any robots.txt rule,
including bot-specific rules, wildcard rules, root blocks, and non-root
path blocks. \\
\textit{SearchBotWild} &
Search-bot wildcard-only block &
Firm blocks at least one search-oriented crawler only through
\texttt{User-agent: *}, with no search-bot-specific rule. \\
\textit{RobotsWild} &
Broad wildcard-only block &
Firm blocks at least one tracked crawler only through wildcard rules,
with no bot-specific rule across the tracked crawler set. \\
\textit{RootSpec} &
Full-site bot-specific block &
Firm adopts a bot-specific rule with \texttt{Disallow: /}, restricting
the full site for the targeted crawler. \\
\bottomrule
\end{tabular}
\TableNote{Treated firms must exhibit an observed clean-to-block transition
within the robots.txt panel. Left-censored and future-treated firms are
excluded. Clean-never firms are retained as controls. Wildcard-only means
that the relevant blocking rule is inherited from \texttt{User-agent: *},
not from a bot-specific group.}
\end{threeparttable}
\end{table}

\subsection{Filing-Window Illiquidity}
\label{subsec:robots_amihud}

I measure filing-window illiquidity using Amihud's price-impact measure,
following \citet{Amihud2002}.  For stock \(i\) on trading day \(d\),
\[
    \text{ILLIQ}_{i,d}
    =
    10^6 \times
    \frac{|R_{i,d}|}{\text{DollarVolume}_{i,d}},
\]
where \(R_{i,d}\) is the daily return and dollar volume is price times
share volume.  For filing event \(e\), let \(W\) denote an event window
such as \([0,+5]\), and let \(B=[-30,-6]\) denote the main pre-filing
benchmark window.  The abnormal Amihud outcome is
\[
    Y_{i,e}(W)
    =
    \frac{1}{|D_{i,e}(W)|}
    \sum_{d\in D_{i,e}(W)}
    \text{ILLIQ}_{i,d}
    -
    \frac{1}{|D_{i,e}(B)|}
    \sum_{d\in D_{i,e}(B)}
    \text{ILLIQ}_{i,d},
\]
where \(D_{i,e}(W)\) is the set of trading days for firm \(i\) in event
window \(W\).  The main event windows are \([0,+1]\), \([0,+3]\),
\([0,+5]\), \([-5,+5]\), \([+6,+30]\), and \([0,+30]\).  The estimating
equation absorbs firm and filing-month fixed effects.  I report both a
two-way fixed effects companion estimate and an imputation DID estimate
following \citet{BorusyakJaravelSpiess2024}.

\begin{table}[!htbp]
\centering
\caption{Filing-Window Amihud Illiquidity: Main \([0,+5]\) Window}
\label{tab:robots_amihud_main}
\begin{threeparttable}
\scriptsize
\resizebox{\textwidth}{!}{%
\begin{tabular}{lrrrrrrrrrr}
\toprule
Treatment &
\multicolumn{4}{c}{BJS imputation DID} &
\multicolumn{6}{c}{TWFE companion} \\
\cmidrule(lr){2-5}\cmidrule(lr){6-11}
& Est. & \(p\) & Pretrend \(p\) & \(N\)
& Est. & \(p\) & \(N\) & Firms & FE & Within \(R^2\) \\
\midrule
SearchBotSpec & 0.054 & 0.006 & 0.535 & 68,289
              & 0.057 & 0.049 & 69,529 & 4,351 & Firm, month & 0.0001 \\
BotSpec       & 0.045 & 0.001 & 0.416 & 66,387
              & 0.046 & 0.033 & 67,693 & 4,242 & Firm, month & 0.0001 \\
RobotsAny     & 0.025 & 0.110 & 0.985 & 29,178
              & 0.019 & 0.213 & 30,881 & 1,947 & Firm, month & 0.0001 \\
SearchBotWild & 0.023 & 0.155 & 0.946 & 29,762
              & 0.016 & 0.336 & 31,411 & 1,979 & Firm, month & 0.0001 \\
RobotsWild    & 0.022 & 0.148 & 0.969 & 31,225
              & 0.015 & 0.350 & 32,872 & 2,067 & Firm, month & 0.0000 \\
RootSpec      & 0.050 & 0.007 & 0.152 & 67,973
              & 0.050 & 0.096 & 69,224 & 4,333 & Firm, month & 0.0000 \\
\bottomrule
\end{tabular}}
\TableNote{The dependent variable is abnormal average Amihud illiquidity
in the \([0,+5]\) filing window. BJS estimates use imputation DID with
firm clustering and 12-month pretrend tests. TWFE estimates absorb firm
and filing-month fixed effects with firm-clustered standard errors; the
last column reports the corresponding within-\(R^2\).}
\end{threeparttable}
\end{table}

Table~\ref{tab:robots_amihud_main} focuses on \([0,+5]\), the main
filing-window horizon.  The search-bot-specific treatment produces a
positive and statistically significant increase in illiquidity.  The
broader \textit{RobotsAny} treatment is attenuated, and wildcard-only
restrictions are not statistically significant in the main window.  This
pattern is consistent with the interpretation that the strongest effect
comes from explicit bot-specific access restrictions rather than generic
site-wide crawler policies.

Table~\ref{tab:robots_amihud_windows} reports the event-window pattern
for the two bot-specific treatments using the TWFE companion estimator,
which provides a conventional within-\(R^2\).  The estimates are
concentrated in the filing window, while the delayed windows
\([+6,+30]\) and \([0,+30]\) are null.

\begin{table}[!htbp]
\centering
\caption{TWFE Event-Window Dynamics for Bot-Specific Robots.txt Treatments}
\label{tab:robots_amihud_windows}
\begin{threeparttable}
\scriptsize
\resizebox{\textwidth}{!}{%
\begin{tabular}{llrrrrrr}
\toprule
Treatment & Window & TWFE Est. & \(p\) & \(N\) & Firms & Fixed effects & Within \(R^2\) \\
\midrule
SearchBotSpec & \([0,+1]\)   &  0.076 & 0.021 & 69,296 & 4,351 & Firm, month & 0.0001 \\
SearchBotSpec & \([0,+3]\)   &  0.055 & 0.055 & 68,427 & 4,351 & Firm, month & 0.0001 \\
SearchBotSpec & \([0,+5]\)   &  0.057 & 0.049 & 69,529 & 4,351 & Firm, month & 0.0001 \\
SearchBotSpec & \([-5,+5]\)  &  0.048 & 0.057 & 69,559 & 4,351 & Firm, month & 0.0001 \\
SearchBotSpec & \([+6,+30]\) & -0.032 & 0.234 & 69,300 & 4,351 & Firm, month & 0.0000 \\
SearchBotSpec & \([0,+30]\)  & -0.017 & 0.447 & 69,385 & 4,351 & Firm, month & 0.0000 \\
\midrule
BotSpec & \([0,+1]\)   &  0.061 & 0.010 & 67,464 & 4,242 & Firm, month & 0.0001 \\
BotSpec & \([0,+3]\)   &  0.044 & 0.037 & 66,629 & 4,242 & Firm, month & 0.0001 \\
BotSpec & \([0,+5]\)   &  0.046 & 0.033 & 67,693 & 4,242 & Firm, month & 0.0001 \\
BotSpec & \([-5,+5]\)  &  0.038 & 0.041 & 67,723 & 4,242 & Firm, month & 0.0001 \\
BotSpec & \([+6,+30]\) & -0.019 & 0.330 & 67,464 & 4,242 & Firm, month & 0.0000 \\
BotSpec & \([0,+30]\)  & -0.008 & 0.647 & 67,549 & 4,242 & Firm, month & 0.0000 \\
\bottomrule
\end{tabular}}
\TableNote{This table reports TWFE companion estimates so that each window
has a conventional within-\(R^2\). The specification absorbs firm and
filing-month fixed effects and clusters standard errors by firm. The BJS
imputation estimates in the main \([0,+5]\) window are reported in
Table~\ref{tab:robots_amihud_main}; the BJS all-window pattern is
qualitatively the same.}
\end{threeparttable}
\end{table}

\begin{table}[!htbp]
\centering
\caption{Pre-Filing Benchmark Robustness for Abnormal Amihud}
\label{tab:amihud_benchmark_robustness}
\begin{threeparttable}
\resizebox{\textwidth}{!}{%
\begin{tabular}{lrrrrrrrrr}
\toprule
Pre-filing benchmark &
\([0,+5]\) Est. & \(p\) & \(N\) & \(R^2\) &
\([+6,+30]\) Est. & \(p\) & \(N\) & \(R^2\) &
Fixed effects \\
\midrule
\([-30,-6]\)  & 0.057 & 0.049 & 69,529 & 0.0001 & -0.032 & 0.234 & 69,300 & 0.0000 & Firm, month \\
\([-45,-6]\)  & 0.051 & 0.080 & 69,532 & 0.0001 & -0.025 & 0.395 & 69,304 & 0.0000 & Firm, month \\
\([-60,-6]\)  & 0.054 & 0.112 & 69,532 & 0.0001 & -0.024 & 0.403 & 69,304 & 0.0000 & Firm, month \\
\([-60,-16]\) & 0.037 & 0.259 & 69,341 & 0.0000 & -0.036 & 0.255 & 69,111 & 0.0000 & Firm, month \\
\bottomrule
\end{tabular}}
\TableNote{The table varies the pre-filing benchmark used to construct
abnormal Amihud for the \textit{SearchBotSpec} treatment. The first row
uses the main \([-30,-6]\) benchmark; the remaining rows move the benchmark
farther from the filing date. The delayed-window placebo remains
insignificant across all benchmark definitions. \(R^2\) is the within-\(R^2\)
from the corresponding TWFE specification.}
\end{threeparttable}
\end{table}

\subsection{Post-Earnings-Announcement Drift}
\label{subsec:robots_pead}

I also test whether robots.txt restrictions amplify
post-earnings-announcement drift.  This test is motivated by three related
strands of the PEAD literature.  \citet{DellaVignaPollet2009} motivate
investor inattention tests around earnings announcements.  I follow this
logic by using announcement-date windows and post-announcement drift
windows.  \citet{HirshleiferLimTeoh2009} study distraction and delayed
reaction to earnings news using analyst forecast errors and drift windows;
I therefore use paper-specific earnings-surprise measures and ask whether
the slope of post-announcement returns with respect to surprise becomes
steeper after robots.txt treatment.  \citet{Martineau2022} shows that PEAD has largely
declined in recent samples, especially outside microcaps.  I therefore
first run a Martineau-style baseline power check and then test treatment
interactions only in baseline-alive PEAD cells.

\begin{table}[!htbp]
\centering
\caption{How the PEAD Literature Maps to the Drift Tests}
\label{tab:pead_literature_map}
\begin{threeparttable}
\small
\begin{tabular}{p{3.2cm}p{4.8cm}p{6.5cm}}
\toprule
Paper & Relevant idea & Implementation here \\
\midrule
\citet{DellaVignaPollet2009} &
Investor inattention around earnings announcements. &
Use earnings-announcement dates and separate immediate announcement
windows from post-announcement drift windows. \\
\citet{HirshleiferLimTeoh2009} &
Distraction can weaken immediate reaction and amplify delayed drift. &
Estimate whether the surprise-return slope changes after robots.txt
treatment; use analyst forecast error and random-walk SUE surprise ranks. \\
\citet{Martineau2022} &
Modern PEAD is weak outside microcaps. &
Run a baseline PEAD power check by microcap status and test treatment
interactions only in cells where baseline drift remains detectable. \\
\bottomrule
\end{tabular}
\end{threeparttable}
\end{table}

Formally, I estimate
\[
    BHAR_{i,q,w}
    =
    \beta\,\text{Surprise}_{i,q}
    +
    \alpha\,\text{PostTreat}_{i,q}
    +
    \gamma\,
    \text{Surprise}_{i,q}\times \text{PostTreat}_{i,q}
    +
    \mu_i+\tau_q+\varepsilon_{i,q,w}.
\]
The coefficient of interest is \(\gamma\).  A positive \(\gamma\) would
imply that robots.txt restrictions amplify drift by steepening the
relation between earnings surprise and post-announcement abnormal returns.

\begin{table}[!htbp]
\centering
\caption{HLT-Style PEAD Interaction Tests}
\label{tab:robots_pead_hlt}
\begin{threeparttable}
\scriptsize
\resizebox{\textwidth}{!}{%
\begin{tabular}{lrrrrrr}
\toprule
Treatment & \([0,+1]\) & \([+2,+61]\) & \(N\) & Treated firms & Fixed effects & Within \(R^2\) \\
\midrule
SearchBotSpec & \(+0.34\) \((0.717)\) & \(-1.00\) \((0.747)\) & 23,807 & 64  & Firm, month & 0.0126 \\
BotSpec       & \(+0.34\) \((0.660)\) & \(-1.65\) \((0.450)\) & 23,131 & 105 & Firm, month & 0.0111 \\
RobotsAny     & \(-0.48\) \((0.198)\) & \(-0.09\) \((0.934)\) & 11,156 & 388 & Firm, month & 0.0148 \\
Broad discovery & \(+0.04\) \((0.939)\) & \(-2.77^{*}\) \((0.063)\) & 22,327 & 173 & Firm, month & 0.0129 \\
\bottomrule
\end{tabular}}
\TableNote{This table follows the HLT-style analyst forecast error
construction: continuous winsorized forecast error scaled by the
quarter-end price. Each cell reports \(\hat{\gamma}\) multiplied by the
headline \(p90-p10\) forecast-error scale, in percentage points; firm-clustered
\(p\)-values are in parentheses. The prediction of delayed diffusion is
weak immediate-window response and positive drift-window amplification.
\(^{*}\) denotes \(p<0.10\).}
\end{threeparttable}
\end{table}

\begin{table}[!htbp]
\centering
\caption{DVP-Style PEAD Interaction Tests}
\label{tab:robots_pead_dvp}
\begin{threeparttable}
\scriptsize
\resizebox{\textwidth}{!}{%
\begin{tabular}{llrrrrr}
\toprule
Surprise form & Treatment & \([0,+1]\) & \([+2,+61]\) & \(N\) & Treated firms & Within \(R^2\) \\
\midrule
Continuous FE & SearchBotSpec & \(+3.47\) \((0.006)\) & \(-1.72\) \((0.370)\) & 18,127 & 52  & 0.0123 \\
Continuous FE & BotSpec       & \(+1.61\) \((0.144)\) & \(-2.27\) \((0.172)\) & 17,552 & 82  & 0.0113 \\
Continuous FE & RobotsAny     & \(-0.60\) \((0.164)\) & \(-1.80\) \((0.108)\) & 8,750  & 321 & 0.0112 \\
Continuous FE & Broad discovery & \(+1.40\) \((0.021)\) & \(-0.33\) \((0.772)\) & 16,891 & 156 & 0.0120 \\
\midrule
Sign-split FE & SearchBotSpec & \(+1.93\) \((0.560)\) & \(-3.75\) \((0.404)\) & 18,127 & 52  & 0.0107 \\
Sign-split FE & BotSpec       & \(-1.31\) \((0.583)\) & \(-2.93\) \((0.387)\) & 17,552 & 82  & 0.0100 \\
Sign-split FE & RobotsAny     & \(-0.91\) \((0.426)\) & \(-1.22\) \((0.628)\) & 8,750  & 321 & 0.0102 \\
Sign-split FE & Broad discovery & \(+3.13^{*}\) \((0.071)\) & \(-4.51^{*}\) \((0.056)\) & 16,891 & 156 & 0.0106 \\
\bottomrule
\end{tabular}}
\TableNote{This table follows the DVP-style construction: forecasts are
drawn from the 30 calendar days before the announcement and scaled by the
price five trading days before the announcement. The continuous rows use
the winsorized forecast error; the sign-split rows use the DVP eleven-bin
forecast-error rank. Reported effects are implied percentage-point changes
in the high-minus-low surprise response; firm-clustered \(p\)-values are in
parentheses. All specifications absorb firm and announcement-month fixed
effects. \(^{*}\) denotes \(p<0.10\).}
\end{threeparttable}
\end{table}

\begin{table}[!htbp]
\centering
\caption{Martineau-Style Baseline-Alive PEAD Cells}
\label{tab:robots_pead_martineau}
\begin{threeparttable}
\scriptsize
\resizebox{\textwidth}{!}{%
\begin{tabular}{lrrrrrrrr}
\toprule
Treatment
& SUE micro \([+2,+5]\)
& SUE micro \([+6,+15]\)
& SUE micro \([+2,+60]\)
& FE full \([+2,+60]\)
& FE micro \([+2,+60]\)
& \(N\) & Treated firms & \(R^2\) \\
\midrule
SearchBotSpec & \(+5.26\) \((0.360)\) & \(+4.52\) \((0.125)\) & \(+0.69\) \((0.887)\) & \(-1.98\) \((0.712)\) & \(+1.93\) \((0.823)\) & 31,418 & 93  & 0.0007 \\
BotSpec       & \(+5.41\) \((0.287)\) & \(+3.65\) \((0.172)\) & \(-1.28\) \((0.785)\) & \(-2.08\) \((0.607)\) & \(-0.87\) \((0.907)\) & 30,734 & 116 & 0.0007 \\
RobotsAny     & \(+1.44^{*}\) \((0.059)\) & \(+0.67\) \((0.465)\) & \(+2.83\) \((0.248)\) & \(-0.41\) \((0.840)\) & \(-0.70\) \((0.842)\) & 14,456 & 416 & 0.0006 \\
Broad discovery & \(+4.49\) \((0.186)\) & \(-0.11\) \((0.942)\) & \(-1.39\) \((0.706)\) & \(-2.82\) \((0.226)\) & \(-2.57\) \((0.655)\) & 31,029 & 144 & 0.0007 \\
\bottomrule
\end{tabular}}
\TableNote{This table follows the Martineau-style design. Surprise ranks
are lagged-quarter empirical deciles; the displayed cells are the
baseline-alive PEAD cells selected by the Martineau power check. Each cell
reports \(\hat{\gamma}\times 9\), the implied \(D10-D1\) drift-spread change
after treatment, in percentage points. Firm-clustered \(p\)-values are in
parentheses. All specifications absorb firm and announcement-quarter fixed
effects. \(^{*}\) denotes \(p<0.10\); no cell reaches \(p<0.05\).}
\end{threeparttable}
\end{table}

The drift results do not support a delayed-earnings-news channel.  Even in
the microcap cells where baseline PEAD remains alive, robots.txt
restrictions do not robustly amplify post-announcement drift.  The only
marginal estimate appears for the broadest \textit{RobotsAny} treatment in
the short \([+2,+5]\) window, and it does not persist over the cumulative
\([+2,+60]\) drift window.  For the broad discovery treatment that produces
the filing-window Amihud effect, the drift-window interactions are null or
negative.  These findings suggest that the main effect operates through a
filing-window liquidity or price-impact channel, rather than through
delayed earnings-news diffusion.

\newpage
\begin{thebibliography}{99}

\bibitem[Amihud(2002)]{Amihud2002}
Amihud, Yakov. 2002.
Illiquidity and stock returns: Cross-section and time-series effects.
\textit{Journal of Financial Markets} 5(1): 31--56.

\bibitem[Borusyak et~al.(2024)]{BorusyakJaravelSpiess2024}
Borusyak, Kirill, Xavier Jaravel, and Jann Spiess. 2024.
Revisiting event-study designs: Robust and efficient estimation.
\textit{Review of Economic Studies} 91(6): 3253--3285.

\bibitem[DellaVigna and Pollet(2009)]{DellaVignaPollet2009}
DellaVigna, Stefano, and Joshua M. Pollet. 2009.
Investor inattention and Friday earnings announcements.
\textit{Journal of Finance} 64(2): 709--749.

\bibitem[Hirshleifer et~al.(2009)]{HirshleiferLimTeoh2009}
Hirshleifer, David, Sonya Seongyeon Lim, and Siew Hong Teoh. 2009.
Driven to distraction: Extraneous events and underreaction to earnings news.
\textit{Journal of Finance} 64(5): 2289--2325.

\bibitem[Martineau(2022)]{Martineau2022}
Martineau, Charles. 2022.
Rest in peace post-earnings announcement drift.
\textit{Critical Finance Review} 11(3--4): 613--646.

\end{thebibliography}

\end{document}
